\section{Local components and exact transport on the original automorphic part}
\label{sec:local-transport-v44}

\status{This chapter proves the seven original local-transport
suppliers \texttt{partUR2IsoL}, \texttt{partUR2IsoQ},
\texttt{partUR2Sph}, \texttt{partUR2Unr}, \texttt{partUR2Wt},
\texttt{partUR2V0}, and \texttt{partUR2Type\_alt}.  It also constructs
an admissible local component at every split place, with the exact
unramified--spherical equivalence, on the original automorphic
function space.  The scalarization and the integral type are the
ones already reconstructed in ADP and ITA.  No Galois attachment,
base-change theorem, or completed \texttt{PartAltAlt} record is
asserted in this chapter.}

\subsection{Objects, conventions, and reusable elementary arguments}

Fix all the original data
\[
 \ell,\ L,\ \Omega,\ s,\ n\ge1,\ V,\ Q,\ c,\ \iota,\ m,\ \tau_0,\
       \lambda,\ d,\ E,\ \xi.
\]
Thus \(L\) is the original CM field, \(\Omega/L\) its specified
algebraic closure, \(s\) the original Hida setting, \(V\) its original
level datum, and \(Q\) its auxiliary-prime datum.  The coefficient DVR
is \(O=c.O\), and the fixed field isomorphism is
\(\iota:\overline{\mathbf Q}_\ell\longrightarrow\mathbf C\).
The integral representation
\(\tau_0:K_T\longrightarrow\operatorname{GL}_m(O)\) is arbitrary
unless a type hypothesis is expressly imposed.  No open-kernel
hypothesis is added to the first six endpoints.

Put \(H=\mathrm{hset}(Q,\tau_0)\),
\(G=U_{\mathbb A}(L,n)\), \(G_f=U_f(L,n)\), and
\(G_\infty=U_\infty(L,n)\).  The group is the identity-form unitary
group of FNL2.  Right translation has the convention
\[
              (\mathscr R_x\phi)(g)=\phi(gx),
              \qquad \mathscr R_x\mathscr R_y=\mathscr R_{xy}.
\]
Let \(I=\operatorname{Emb}_\ell((s.w(v))_{v\in V.\mathrm{frame}.V_\ell})\),
and use exactly the prescribed basis equivalence
\[
 D=\bigotimes_{a\in I}\nabla_{\lambda(a.\mathrm{res})}(\mathbf C),
       \qquad E:\mathbf C^d\xrightarrow{\sim}D .
\]
The source's \(\mathrm{IsBasis}(\iota,E,\xi)\) is assumed.  Write
\(R=\mathrm{infRep}(\iota,\lambda)\) and retain the original projection
\[
 p(g)=\bigl(\mathrm{away}(g),
                 (\mathrm{cLoc}_{s.w(v)}(g))_{v\in V.\mathrm{frame}.V_\ell}\bigr)
             \in G^a\times G^\ell .
\]
For a coefficient function \(f:G^a\times G^\ell\to\mathbf C^d\),
its scalarization is
\begin{equation}\label{eq:ltp-scalarization}
       \mathcal T_f(\alpha)(g)
       =\alpha\bigl(R(g^{-1})E(f(p(g)))\bigr),
                   \qquad\alpha\in D^* .
\end{equation}
For \(F=(F_i)_{i<m}\), use this formula on each \(F_i\).
Let \(\sigma\) be an original \(\mathrm{IsAutRepU}\) subspace of
the actual functions \(G\to\mathbf C\), and set
\(Y=\mathrm{YOf}(Q,\iota,\tau_0,E,\xi,\sigma)\).

The inherited statements used here have written proofs: FNL2--FNL6
and FNL16 for the actual groups and sections; ADP1--ADP3 and
ADP6--ADP11 for the positive form, finite multiplicity and
scalarization; HP1--HP4 for the characteristic-zero polynomial
representations; and ITA9--ITA12 for the supplied integral type.
The local-transport arguments introduce no new local Langlands,
base-change or Galois-attachment premise.  When ITA12 is used, its
expressly recorded classical type and correspondence foundations
remain at exactly their previous scope.  In particular we do not
turn a supplied \(\mathrm{IsTypeAlt}\) into an unproved assertion
about all local representations.

% BEGIN LTP1
\paragraph{LTP1. The literal part and its finite separating scalarizations.}
\label{claim:LTP1}
For every coefficient vector \(F\),
\[
 F\in Y
 \quad\Longleftrightarrow\quad
 F\in H.\mathrm{Big}(\iota,\xi)
       \ \hbox{and}\ 
 \mathcal T_{F_i}(\alpha)\in\sigma
                  \quad\hbox{for every }i,\alpha .
\]
If \(\alpha_1,\ldots,\alpha_d\) is a basis of \(D^*\), the finitely
many linear maps
\(F\mapsto\mathcal T_{F_i}(\alpha_b)\), indexed by \(i<m,b\le d\),
separate coefficient vectors.  These assertions include \(m=0\)
and \(d=0\).
\begin{proof}
Scalarization is complex linear in \(F_i\) and in \(\alpha\).
Consequently the set of vectors for which all scalarizations belong
to the submodule \(\sigma\) is itself a submodule.  Its span, which
occurs in the literal definition of \(\mathrm{YOf}\), is therefore
that same set.  Intersecting with \(H.\mathrm{Big}\) proves the
equivalence, rather than just one implication for generators.

If all displayed coordinates vanish, linearity in \(\alpha\) makes
every scalarization vanish.  ADP7 supplies injectivity on the original
functions: \(p|_{G_f}\) is surjective, \(R(g)\) is invertible, \(E\)
is an isomorphism, and the full dual separates \(D\).  Thus every
\(F_i=0\).  When an index set is empty the corresponding coefficient
space is zero, so the same assertion holds.  This proves the full
source \texttt{partUR2Mem} and the finite separating family needed
below.
\end{proof}
% END LTP1

% BEGIN LTP2
\paragraph{LTP2. Actual local factors, centralizers, and the opposite place.}
\label{claim:LTP2}
For a split finite place \(w\) of \(L\), there is a continuous
homomorphism
\[
       j_w:\operatorname{GL}_n(L_w)\longrightarrow G_f
\]
whose image is exactly \(\mathrm{locU}(n,w)\), with
\(\mathrm{cLoc}_w j_w(g)=g\).  It has the actual finite coordinates
specified in FNL4--FNL6.  Every \(x\in G\) factors as
\[
       x=j_w(\mathrm{cLoc}_w x)\,x^{(w)},
             \qquad x^{(w)}j_w(g)=j_w(g)x^{(w)}
                                     \quad\hbox{for every }g.
\]
The factor \(x^{(w)}\) has identity coordinates at both places above
\(w|_{L^+}\).  Distinct base places give commuting local factors.
If \(w'\) has the same base place as \(w\), then
\[
 \mathrm{cLoc}_w x\in\operatorname{GL}_n(\mathcal O_{L_w})
 \quad\Longrightarrow\quad
 \mathrm{cLoc}_{w'}x\in\operatorname{GL}_n(\mathcal O_{L_{w'}}).
\]
This last implication does not require that the base place split.
\begin{proof}
At a split pair \(w,w^c\), set the \(w\)-coordinate equal to \(g\),
the other equal to its inverse conjugate transpose, and every
remaining finite and infinite coordinate equal to one.  The unitary
identity verifies membership.  The inverse-transpose operation is
a group homomorphism in this order:
\((gh)^{-t}=g^{-t}h^{-t}\).  Thus the coordinate construction is a
homomorphism.  It is continuous because the two nonconstant
coordinate maps are continuous.  Projection at \(w\) is its inverse
on \(\mathrm{locU}(n,w)\): the local-support condition removes all
other base places and the unitary equation determines the opposite
coordinate.  This is precisely the section already checked in FNL.

Removing these two coordinates from \(x\) gives \(x^{(w)}\).
Coordinate multiplication proves the displayed factorization and
commutation.  It also proves commutation for sections supported at
distinct base places.  Thus the centralizer assertion is about the
actual full adelic group, including its infinite coordinates.

In a quadratic CM extension, two places over one finite base place
are equal or are interchanged by conjugation.  In the second case
the unitary equation gives
\[
 \mathrm{cLoc}_{w'}x
       ={}^t\!\bigl(c_{w,w'}(\mathrm{cLoc}_w x)^{-1}\bigr).
\]
Local conjugation carries the valuation ring onto the valuation
ring.  Membership in \(\operatorname{GL}_n(\mathcal O)\) says that
both a matrix and its inverse have integral entries.  The displayed
formula exchanges those two conditions and preserves them.
This proves the last assertion, including the inert case when the
places coincide.  The source's \(KGL\) and \(KO\) both express this
same lattice-preservation condition.
\end{proof}
% END LTP2

% BEGIN LTP3
\paragraph{LTP3. A local component makes the whole restriction isotypic.}
\label{claim:LTP3}
Suppose
\[
 \mathrm{IsCompAlong}\!
 \left(\sigma,\mathrm{locU}(n,w),\mathrm{cLoc}_w,\tau\right)
\]
for a smooth irreducible \(\tau\) at a split place \(w\).
Then the restriction of the full \(\sigma\) along \(j_w\) is
\(\tau\)-isotypic in the literal \(\mathrm{IsoSum}\) sense.
\begin{proof}
The component hypothesis supplies a nonzero equivariant linear map
from \(\tau\) into the original function subspace \(\sigma\).  It
intertwines the action along \(j_w\), by LTP2.  Let \(Z\) be the sum
of the images of all such local intertwiners.  This sum is nonzero.
Each image is stable under \(j_w(\operatorname{GL}_n(L_w))\).
An element \(h\) centralizing that local subgroup sends an intertwiner
\(f\) to the intertwiner \(\mathscr R_hf\).  Therefore it preserves
\(Z\).  The factorization in LTP2 now shows that every \(x\in G\)
preserves \(Z\).  The original global irreducibility of \(\sigma\)
gives \(Z=\sigma\).  This equality is precisely \(\mathrm{IsoSum}\);
no irreducibility of the whole restriction is asserted.
\end{proof}
% END LTP3

% BEGIN LTP4
\paragraph{LTP4. Finite equivariant separating families preserve isotypy.}
\label{claim:LTP4}
Let \(Y\) be a stable subspace of an arbitrary representation
\(\rho\) of a group \(J\).  Suppose \(Z\) is a \(\tau\)-isotypic
subspace of another representation, and finitely many linear maps
\(\phi_a\) take \(Y\) into \(Z\), intertwine the \(J\)-actions on
\(Y\), and jointly have zero kernel on \(Y\).  Then \(Y\) is
\(\tau\)-isotypic.
\begin{proof}
The combined map embeds \(Y\) equivariantly in \(Z^A\).
An isotypic module, defined as the sum of images of maps from a
simple module \(\tau\), is a sum of simple submodules of that type.
The elementary module argument of ITA10 makes it a direct sum of
copies of \(\tau\), and every submodule of that sum is again a sum
of copies of \(\tau\).  Since \(A\) is finite, \(Z^A\) is a finite
direct sum of such modules, hence has that same property.  Apply
it to the image of \(Y\), and transport the embeddings of the
simple summands back through the injective map.
All statements are algebraic; neither admissibility nor finite
dimension of \(Y,Z,\tau\) is used.
\end{proof}
% END LTP4

% BEGIN LTP5
\paragraph{LTP5. Compact fixed vectors and the exact spherical sums.}
\label{claim:LTP5}
Let \(Z\) be a \(\tau\)-isotypic subspace of a representation
\(\rho\) of \(\operatorname{GL}_n(F)\), where \(F\) is a
nonarchimedean local field and \(\tau\) is smooth irreducible.
For every compact subgroup \(K\),
\[
               Z^K\ne0\quad\Longrightarrow\quad\tau^K\ne0 .
\]
Put \(K_0=\operatorname{GL}_n(\mathcal O_F)\).  If the source
operator \(\mathrm{sphOp}_{\tau}(j)\) equals \(a\) on \(\tau^{K_0}\),
then \(\mathrm{sphOp}_{\rho}(j)\) equals \(a\) on \(Z^{K_0}\),
for every \(j\in\mathbf N\).  Every additive equivariant map
commutes with the source spherical operator.  Finally each original
\(K_c(F,n,c)\), for \(c\ge0\), is compact.
\begin{proof}
For a smooth vector \(v\), its \(K\)-orbit is finite: the stabilizer
is open and its cosets cover the compact group.  Averaging over the
finite left cosets of this stabilizer gives a \(K\)-fixed vector
\(e_Kv\).  Refining stabilizers shows that this is a linear operation,
fixes every \(K\)-fixed vector, and commutes with equivariant linear
maps.  Normality of a stabilizer is unnecessary: left multiplication
permutes its left cosets.

Write \(z\in Z^K\) as a finite sum \(\sum_b f_b(v_b)\) of images of
intertwiners from \(\tau\).  Applying the average gives
\(z=\sum_bf_b(e_Kv_b)\).  If \(\tau^K=0\), this makes \(z=0\).
This proves the first assertion and gives the decomposition into
images of fixed vectors needed for the second.

The original spherical operator is the unnormalized double-coset
sum for
\[
 a_j=\operatorname{diag}
       (\underbrace{\varpi,\ldots,\varpi}_{\min(j,n)},
        1,\ldots,1).
\]
For the source's finite right-coset representatives
\(K_0a_jK_0=\coprod_b g_bK_0\), it is
\(\sum_b\rho(g_b)\).  Finiteness follows because
\(K_0\cap a_jK_0a_j^{-1}\) is open in compact \(K_0\);
this is the same finite-sum dictionary proved in OHA1.
An additive intertwiner commutes with this finite sum term by
term, even before one restricts to fixed vectors.  On the
decomposition \(z=\sum_bf_b(v_b)\) with \(v_b\in\tau^{K_0}\),
the assumed scalar equation therefore gives
\(\mathrm{sphOp}_\rho(j)z=az\).
The argument covers \(j=0\) and \(j\ge n\), with no change of
normalization or factor order.

The subgroup \(K_0\) is compact: identify it with the closed set of
pairs \((A,B)\) in two finite matrix powers of the compact ring
\(\mathcal O_F\) satisfying \(AB=BA=1\).
The original \(K_c(F,n,c)\) is the kernel of reduction modulo
\(\mathfrak m_F^c\), inside this group.  Reduction is continuous
with finite target, so the kernel is closed and compact.
For \(c=0\) the reduction condition is vacuous and the group is
\(K_0\).  These are all four clauses of the source
\texttt{partUR2Abs}.
\end{proof}
% END LTP5

\subsection{Transport of the two local actions}

% BEGIN LTP6
\paragraph{LTP6. The exact action at a place above \(\ell\).}
\label{claim:LTP6}
For \(v\in V.\mathrm{frame}.V_\ell\), use the actual split section
\(j_v=j_{s.w(v)}\).  For every coefficient vector \(F\), every
\(g\in\operatorname{GL}_n(L_{s.w(v)})\), every \(i,\alpha\),
\[
 \mathcal T_{(H.\mathrm{atV}(v,g)F)_i}(\alpha)
                 =\mathscr R_{j_v(g)}\mathcal T_{F_i}(\alpha).
\]
Both \(H.\mathrm{Big}(\iota,\xi)\) and \(Y\) are stable under this
action.
\begin{proof}
The chosen place lies above \(v\), and the splitting condition in
the Hida setting gives LTP2's section.  Its away coordinates are
one and its selected \(\ell\)-coordinates are one except at \(v\),
where the coordinate is \(g\).  Thus
\[
             p(xj_v(g))=p(x)(1,\delta_v(g)).
\]
Its infinite coordinates are one, hence \(R(j_v(g))=1\) by ADP6.
In ADP7's right-translation rule the functional
\(\alpha\circ R(j_v(g)^{-1})\) is therefore exactly \(\alpha\).
The result is the displayed identity, with \(xj_v(g)\) in that
order.  This proves the full source \texttt{partUR2AtV}.

The local right translation commutes with rational left
translation.  It also commutes with the entire away action,
including its coefficient matrix, since the two group coordinates
are distinct and the matrix is constant.  Therefore it preserves
both defining clauses of \(H.\mathrm{Big}\).  Finally the
scalarization identity and global stability of \(\sigma\) preserve
all the scalarizations required by LTP1, proving stability of \(Y\).
\end{proof}
% END LTP6

% BEGIN LTP7
\paragraph{LTP7. Full original \texttt{partUR2IsoL}.}
\label{claim:LTP7}
For every \(v\in V.\mathrm{frame}.V_\ell\) and every supplied smooth
irreducible local component \(\tau\) of \(\sigma\) along
\(\mathrm{locU}(n,s.w(v))\),
\[
                  \mathrm{IsoSum}(H.\mathrm{atV}(v),Y,\tau).
\]
No nonzero-\(Y\) or open-\(\ker\tau_0\) assumption is required.
\begin{proof}
LTP3 makes the restriction of \(\sigma\) along \(j_v\)
\(\tau\)-isotypic.  LTP6 gives the stability and intertwining
conditions for the finite scalarization family of LTP1.
That family is separating and lands in \(\sigma\).
LTP4 proves the stated literal \(\mathrm{IsoSum}\).
\end{proof}
% END LTP7

% BEGIN LTP8
\paragraph{LTP8. The exact away action, including its coefficient matrix.}
\label{claim:LTP8}
Let \(w\) be split with \(w|_{L^+}\notin s.S\).  For every
coefficient vector \(F\), \(g\in\operatorname{GL}_n(L_w)\),
and \(i,\alpha\),
\[
 \mathcal T_{(H.\mathrm{atW}(w,g)F)_i}(\alpha)
                    =\mathscr R_{j_w(g)}\mathcal T_{F_i}(\alpha).
\]
At an original auxiliary place \(w=Q.wq(q)\), the actual
\(H.\mathrm{Ax}\) makes \(H.\mathrm{Big}\) and \(Y\) stable
under this action.
\begin{proof}
The place lies away from \(\ell\), the type places, and \(v_0\).
The away projection of \(j_w(g)\) belongs to the original
\(\mathrm{locA}(n,s,w)\).  Its \(w\)-coordinate is \(g\);
all other base coordinates are one.  The bijective coordinate
projection checked in FNL4--FNL6 therefore identifies that away
projection with the source's \(\mathrm{sec}(w,g)\), including
the branch of its conditional definition.  Furthermore
\[
 \mathrm{toKT}(\mathrm{sec}(w,g))=1,\qquad
 H.\tau_C(\mathrm{sec}(w,g))=1.
\]
Its selected \(\ell\)-coordinates and all infinite coordinates
are one.  Thus the original matrix-valued away action
\[
  (\mathrm{raF}(u)F)_i(y)
       =\sum_b\tau_C(u)_{ib}F_b(y(u,1))
\]
reduces at \(u=\mathrm{sec}(w,g)\) to the pure right translation
\(F_i(y(u,1))\).  ADP7 gives exactly the asserted scalarization
identity.  This is the full source \texttt{partUR2AtW}.

For an auxiliary \(q\), the original axiom \(\mathrm{Q\_UQ}\)
says that \(\mathrm{sec}(Q.wq(q),g)\) commutes with every element
of \(H.UQ\).  Consequently its \(\mathrm{raF}\) action preserves
the \(UQ\)-fixed condition.  It also preserves rational coefficient
equivariance: every coordinate is a finite complex linear
combination of right translates, which commute with rational
left translation.  It preserves \(H.\mathrm{Big}\).
The scalarization identity and LTP1 then prove stability of \(Y\).
The axiom has been verified for this very \(H\) in FNL16.
\end{proof}
% END LTP8

% BEGIN LTP9
\paragraph{LTP9. Full original \texttt{partUR2IsoQ}.}
\label{claim:LTP9}
Under the original \(H.\mathrm{Ax}\), for every auxiliary
\(q\in Q.Q\) and every supplied smooth irreducible local component
\(\tau\) of \(\sigma\) at \(Q.wq(q)\),
\[
             \mathrm{IsoSum}(H.\mathrm{atW}(Q.wq(q)),Y,\tau).
\]
This includes \(Y=0\).
\begin{proof}
The auxiliary-prime conditions give
\(Q.wq(q)|_{L^+}\notin s.S\) and membership of this base place
in \(V.S'\); the original level datum says it splits.
Thus LTP2, LTP3 and LTP8 apply.  The same finite family from
LTP1 satisfies all the hypotheses of LTP4, giving the conclusion.
The original \(H.\mathrm{Ax}\) premise is retained in this generic
endpoint and is available from FNL16 in the actual application.
\end{proof}
% END LTP9

\subsection{The compact level equation and its three consequences}

% BEGIN LTP10
\paragraph{LTP10. Full matrix-valued transport of the base level equation.}
\label{claim:LTP10}
Let \(F\in H.\mathrm{Big}(\iota,\xi)\), and let \(x\in G_f\).
Assume its selected \(\ell\)-coordinates are one.  Let \(k\in K_T\)
have components \(k_u=\mathrm{cLoc}_{s.w(u)}x\) at every type
place \(u\in s.St\).  Suppose also that
\[
\begin{gathered}
 \mathrm{cLoc}_{s.w(v_0)}x\in K_c(L_{s.w(v_0)},n,1),\\
 \mathrm{cLoc}_{w}x\in KO_w
                  \quad(w|_{L^+}\notin V.S'),\qquad
 \mathrm{cLoc}_{Q.wq(q)}x=1\quad(q\in Q.Q).
\end{gathered}
\]
Writing \(t=\iota(\tau_0)\) for the original coefficient extension,
one has, for every \(i,\alpha\),
\begin{equation}\label{eq:ltp-level-matrix}
       \sum_{b<m} t(k)_{ib}\,
            \mathscr R_x\mathcal T_{F_b}(\alpha)
                      =\mathcal T_{F_i}(\alpha).
\end{equation}
This is the full source \texttt{partUR2Fix}.
\begin{proof}
Set \(u=\mathrm{away}(x)\).  The definition of \(H.U_b\) asks that
its type components be integral and its distinguished component
lie in \(K_c(1)\); these follow from the stated conditions.
The definition of \(H.UQ\) adds precisely integrality outside
\(V.S'\) and the identity at each auxiliary place; these are the
remaining displayed conditions.  The source's away-coordinate
map \(\mathrm{cmp}\) equals \(\mathrm{cLoc}\) away from \(\ell\)
and equals one at a place above \(\ell\), so all these subgroup
membership assertions are valid on the literal away projection.
The type places lie away from \(\ell\), by \(s.St_\ell\).
Consequently \(\mathrm{toKT}(u)=k\) exactly.

The \(UQ\)-fixed clause of \(H.\mathrm{Big}\) is therefore the
coordinate equation
\[
       \sum_b t(k)_{ib}F_b(y(u,1))=F_i(y)
                     \quad\hbox{for every }y.
\]
Since \(x\) is finite, \(R(x)=1\); since its selected
\(\ell\)-coordinates are one, \(p(gx)=p(g)(u,1)\).
Apply \(E\), then \(R(g^{-1})\), and finally \(\alpha\) to the
last equation at \(y=p(g)\).  Linearity gives
\eqref{eq:ltp-level-matrix}.  In particular the coefficient
matrix has not been discarded or replaced by a scalar character.
\end{proof}
% END LTP10

% BEGIN LTP11
\paragraph{LTP11. Full original \texttt{partUR2Sph}.}
\label{claim:LTP11}
Assume the original \(H.\mathrm{Ax}\).
Let \(w|_{L^+}\notin V.S'\) be split, let \(\tau\) be any
supplied local component of \(\sigma\) at \(w\), and let
\(j\in\mathbf N\), \(a\in\mathbf C\).  If
\[
       \mathrm{sphOp}_\tau(j)v=av
                      \quad\hbox{for every }v\in\tau^{KO_w},
\]
then, for every \(F\in Y\),
\[
                \mathrm{sphOp}_{H.\mathrm{atW}(w)}(j)F=aF .
\]
\begin{proof}
Here \(w|_{L^+}\notin s.S\), so LTP8 supplies the exact
intertwining formula on all coefficient vectors.  The axiom
\(\mathrm{W\_le}\) says that the image of \(KO_w\) under the
original section lies in \(H.UQ\).  Since \(F\in H.\mathrm{Big}\),
it follows that \(H.\mathrm{atW}(w,k)F=F\) for every \(k\in KO_w\).

LTP3 makes \(\sigma\), restricted along \(j_w\), \(\tau\)-isotypic.
Each scalarization of \(F\) lies in that subspace and is \(KO_w\)-fixed.
LTP5 transfers the scalar equation to it.  The same claim's
finite-sum naturality, applied to the linear scalarization maps
and LTP8, gives
\[
 \mathcal T_{(\mathrm{sphOp}_{H.\mathrm{atW}(w)}(j)F)_i}(\alpha)
                          =a\,\mathcal T_{F_i}(\alpha).
\]
LTP1 separates the two coefficient vectors, proving the result.
We used the chosen source double-coset sum, not a rescaled
Hecke operator.  No nonzero-\(Y\) assumption, upper bound on \(j\),
or extra condition on \(\tau_0\) has entered.
\end{proof}
% END LTP11

% BEGIN LTP12
\paragraph{LTP12. Full original \texttt{partUR2Unr}, also at nonsplit places.}
\label{claim:LTP12}
If \(Y\ne0\), then for every finite place \(w\) of \(L\) with
\(w|_{L^+}\notin V.S'\), the original
\(\mathrm{IsUnramU}(\sigma,w)\) holds.
\begin{proof}
Choose \(0\ne F\in Y\).  By LTP1 some
\(\phi=\mathcal T_{F_i}(\alpha)\) is nonzero, and \(\phi\in\sigma\).
We show that this one function is fixed by every
\(x\in\mathrm{locU}(n,w)\) whose \(w\)-coordinate belongs to
the original \(KGL_w=KO_w\).

Such \(x\) is finite and is supported above the base place
\(w|_{L^+}\).  Since that place lies outside \(V.S'\), it lies
outside every selected \(\ell\)-place, every type place, \(v_0\),
and every auxiliary place.  All of \(x\)'s coordinates there
are one.  At a different place \(w'\) above \(w|_{L^+}\), the
last assertion of LTP2 makes its coordinate integral; at all
other base places its coordinates are one.  Thus every hypothesis
of LTP10 holds with \(k=1\).
As \(t(1)\) is the identity matrix,
\eqref{eq:ltp-level-matrix} gives \(\mathscr R_x\phi=\phi\).

The chosen \(\phi\) is nonzero, belongs to \(\sigma\), and has
exactly the invariance required in \(\mathrm{IsUnramU}\).
The argument used the local-support subgroup directly and did
not use a split section at \(w\); it therefore proves the
literal endpoint at nonsplit places as well.
\end{proof}
% END LTP12

\subsection{The prescribed dual infinity weight}

% BEGIN LTP13
\paragraph{LTP13. The actual CM embeddings and the single-factor dual map.}
\label{claim:LTP13}
The map
\[
       a\longmapsto\operatorname{InfinitePlace}
                       (\iota\circ a.\mathrm{res}),\qquad a\in I,
\]
is injective.  For a finite family \(D_a\) of finite-dimensional
complex spaces, any nonzero linear map
\(\Phi:(\bigotimes_aD_a)^*\to N\), and every \(a\), there is a
linear map
\[
                 e_a:D_a^*\longrightarrow(\bigotimes_bD_b)^*
\]
such that \(\Phi e_a\ne0\), which intertwines precomposition by
every tensor operator acting as the identity outside factor \(a\).
For the actual polynomial module,
\[
 \nabla_{\mu^\vee}(\mathbf C)
       \simeq\nabla_\mu(\mathbf C)^*,\qquad
              \mu^\vee_i=-\mu_{n-1-i},
\]
as representations with contragredient action on the right.
The displayed \(\mu^\vee\) is exactly the source's
\(\mathrm{DomWeight.dual}\).
\begin{proof}
Write \(a=(v,j)\), where \(j:L_{s.w(v)}\to
\overline{\mathbf Q}_\ell\) is one of the original continuous
local embeddings.  Such an embedding recovers the chosen prime:
for \(z\in\mathcal O_L\),
\[
                z\in s.w(v)
                   \quad\Longleftrightarrow\quad |j(z)|_\ell<1.
\]
This is the finite local norm calculation of WRP1.  It also follows
by applying continuity to powers of \(z\), and using maximality
of the two resulting prime ideals.

Equal infinite places give equality or complex conjugacy of the
two complex embeddings.  In the equality case, injectivity of
\(\iota\) identifies their restrictions to \(L\), hence their
primes and their base places by the displayed test.  The local
embeddings then agree by continuity and density of \(L\) in its
completion.

In the conjugacy case the CM identity
\(\overline{\iota(j(z))}=\iota(j(cz))\), for \(z\in L\),
identifies one restriction with the other composed with CM
conjugation.  It first gives equality on \(L^+\), hence the same
base place \(v\).  Both chosen places must consequently be the
same \(s.w(v)\), while the prime test makes one the conjugate of
the other.  This contradicts splitting at \(v\).  Thus conjugacy
cannot occur between two indices.  No continuity of \(\iota\)
has been used; the CM identity is an identity of embeddings of
the number field.

For the tensor assertion, pure tensors of dual vectors span the
dual tensor product, by finite-dimensional duality.  Choose
\((\beta_b)_b\) with
\(\Phi(\bigotimes_b\beta_b)\ne0\), and define \(e_a(\beta)\)
by replacing only \(\beta_a\) with \(\beta\).
Then \(\Phi e_a(\beta_a)\ne0\).  Evaluation on a pure tensor
shows directly that
\[
 e_a(\beta)\circ(\bigotimes_b T_b)
          =e_a(\beta\circ T_a)
                       \quad\hbox{if }T_b=1\ (b\ne a).
\]

Finally HP1--HP4 identify the exact \(\nabla_\mu\) as the
irreducible rational representation with highest weight \(\mu\).
Its dual is irreducible: the annihilator of a proper nonzero
dual-stable subspace would be a proper nonzero original-stable
subspace.  Its highest weight is the negative of the lowest
weight of the original representation, namely
\((-\mu_{n-1},\ldots,-\mu_0)\).
Indeed the longest permutation carries the highest weight line
to the lowest weight line, and the functional on that line
annihilating the other weight spaces is fixed by the upper
unipotents in the contragredient action.  The weight calculation
and classification of HP4 therefore give the displayed
isomorphism on the original polynomial modules.  This calculation
specifies a single dual at the same complex embedding; it does
not silently change to a conjugate CM type.
\end{proof}
% END LTP13

% BEGIN LTP14
\paragraph{LTP14. Full original \texttt{partUR2Wt}.}
\label{claim:LTP14}
If \(Y\ne0\), then for every \(a\in I\),
\[
   \mathrm{HasWeightAt}
       \bigl(\sigma,\iota\circ a.\mathrm{res},
                         \lambda(a.\mathrm{res})^\vee\bigr).
\]
\begin{proof}
Choose \(F\in Y\) with one coordinate \(F_i\ne0\).  By ADP7
the linear map
\[
                 \Phi:D^*\longrightarrow\sigma,\qquad
                        \alpha\longmapsto\mathcal T_{F_i}(\alpha)
\]
is nonzero.  If \(x\in G_\infty\), then \(p(x)=1\), so the
exact translation rule reads
\[
       \mathscr R_x\Phi(\alpha)
                         =\Phi(\alpha\circ R(x^{-1})).
\]
Fix \(a\).  On the source subgroup
\(\mathrm{embU}(n,\iota\circ a.\mathrm{res})\), every other
tensor factor of \(R\) is the identity, by LTP13's injectivity
of infinite places.  Use its map \(e_a\) and the dual
isomorphism there.  The composite
\[
 \nabla_{\lambda(a.\mathrm{res})^\vee}(\mathbf C)
       \longrightarrow D_a^*
       \xrightarrow{e_a}D^*
       \xrightarrow{\Phi}\sigma
\]
is nonzero and intertwines this exact embedded unitary action,
with the source's \(\mathrm{cEmb}\) projection.
Its domain is the irreducible polynomial representation
required in \(\mathrm{IsCompAlong}\), by HP2.
These are precisely the two clauses of
\(\mathrm{HasWeightAt}\).  The inverse in the translation rule
is responsible for the dual weight; no sign or determinant shift
is omitted.  The assertion is vacuous when \(I\) is empty.
\end{proof}
% END LTP14

% BEGIN LTP15
\paragraph{LTP15. Full original \texttt{partUR2V0} at depth one.}
\label{claim:LTP15}
If \(Y\ne0\), then every supplied smooth irreducible local
component \(\tau\) of \(\sigma\) at the distinguished place
\(s.w(v_0)\) satisfies
\[
                     \tau^{K_c(L_{s.w(v_0)},n,1)}\ne0 .
\]
\begin{proof}
The original \(v_0\) has norm equal to a rational prime \(p\)
different from \(\ell\).  In particular \(p\in v_0\).
If also \(\ell\in v_0\), a Bezout relation for the distinct
rational primes \(p,\ell\), viewed in \(\mathcal O_{L^+}\),
would put \(1\) in \(v_0\).  Thus \(v_0\) is away from \(\ell\).
The original setting also says \(v_0\notin s.St\), and it is
distinct from every auxiliary base place because those places
lie outside \(s.S\).  The setting's splitting condition supplies
the actual section \(j_0=j_{s.w(v_0)}\).

Choose a nonzero scalarization \(\phi\in\sigma\) of a nonzero
\(F\in Y\), as in LTP12.  For
\(g\in K_c(L_{s.w(v_0)},n,1)\), the finite element \(j_0(g)\)
has identity selected \(\ell\)-coordinates, identity type
coordinates, identity auxiliary coordinates, and identity
coordinates outside \(V.S'\).  Its distinguished coordinate is
exactly \(g\), so LTP10 applies with \(k=1\) and gives
\(\mathscr R_{j_0(g)}\phi=\phi\).  Consequently the restriction
of \(\sigma\) has a nonzero vector fixed by this exact
depth-one subgroup.

LTP3 makes that restriction \(\tau\)-isotypic.  The subgroup is
compact by LTP5, so LTP5 transfers nonzero fixedness to \(\tau\).
No shrinking to an unspecified deeper subgroup is made.
\end{proof}
% END LTP15

\subsection{The type equation with every local correspondence retained}

% BEGIN LTP16
\paragraph{LTP16. One actual product section at the type places.}
\label{claim:LTP16}
There is a homomorphism
\[
       J:\prod_{u\in s.St}\operatorname{GL}_n(L_{s.w(u)})
                                      \longrightarrow G_f
\]
such that \(J(\delta_u(g))=j_{s.w(u)}(g)\), its \(s.w(u)\)-coordinate
is \(g_u\), and all coordinates at other base places are one.
In particular it has identity selected \(\ell\)-coordinates,
identity distinguished coordinate, and identity auxiliary coordinates.
For every \(F\in H.\mathrm{Big}\), \(k\in K_T\), and \(i,\alpha\),
\[
       \sum_b t(k)_{ib}\,
               \mathscr R_{J(k)}\mathcal T_{F_b}(\alpha)
                          =\mathcal T_{F_i}(\alpha).
\]
\begin{proof}
Every original type place splits, and the chosen place lies above
its indicated base place, by FNL4.  The base places indexed by
\(s.St\) are distinct.  Their sections therefore commute by LTP2.
Define \(J(g)\) as the finite product of their images.
Commutation makes the product independent of order and makes
\(J\) a homomorphism.  Its coordinates have the stated values
by construction.  This remains valid for an empty type set.

The type places lie away from \(\ell\), exclude \(v_0\), and are
disjoint from the auxiliary base places.  They are all contained
in \(V.S'\).  Consequently \(J(k)\) has precisely the coordinates
required by LTP10: its type coordinates are the prescribed \(k\),
its distinguished coordinate is the identity in \(K_c(1)\), and
outside \(V.S'\) it is the identity.  Apply LTP10 to obtain the
displayed matrix equation.
\end{proof}
% END LTP16

% BEGIN LTP17
\paragraph{LTP17. Full original \texttt{partUR2Type\_alt}.}
\label{claim:LTP17}
Assume \(Y\ne0\).  Let
\(\chi_{u,i}:I_{s.A(u)}\to\mathbf C^\times\) be any family
for which the original
\[
      \mathrm{IsTypeAlt}\bigl((s.hw(u))_u,t,\chi\bigr)
\]
holds.  Let \((\tau_u)_{u\in s.St}\) be any family of smooth
irreducible local components of \(\sigma\) at the chosen type
places.  Then, for every \(u\in s.St\) and every map
\(\mathrm{rc}\) satisfying the original
\(\mathrm{IsRecT}(s.hw(u),\mathrm{rc})\),
\[
        \mathrm{IsTypeWD}\!
          \left(s.hw(u),\chi_u,
                         \mathrm{rc}(\mathrm{algNorm}(\tau_u))\right).
\]
\begin{proof}
Choose \(0\ne F\in Y\).  LTP1 gives \(i_0,\alpha\) such that
\(\mathcal T_{F_{i_0}}(\alpha)\ne0\).  Keep this \(\alpha\)
fixed, and put
\(\Phi_i=\mathcal T_{F_i}(\alpha)\in\sigma\).
Thus \(\Phi=(\Phi_i)_i\) is a single nonzero coefficient vector
in the original nonzero vector space \(\sigma\).

Use LTP16's one homomorphism \(J\) to define a representation
of \(\prod_u\operatorname{GL}_n(L_{s.w(u)})\) on \(\sigma\):
\(\rho(g)\phi=\mathscr R_{J(g)}\phi\).
Global stability of \(\sigma\) makes this well defined.
Restriction to its \(u\)-th factor is exactly the action along
\(j_{s.w(u)}\), so LTP3 says that this restricted representation
is \(\tau_u\)-isotypic.  Inclusion of \(\sigma\) into the ambient
functions identifies this assertion with the literal
\(\mathrm{IsoSum}(\rho|_u,\top,\tau_u)\).

LTP16 proves, for every \(k\in K_T\),
\[
                   \sum_b t(k)_{ib}\rho(k)\Phi_b=\Phi_i .
\]
We have therefore verified the entire left side of the equivalence
in the supplied \(\mathrm{IsTypeAlt}\), on this very nonzero
module \(\sigma\) and this product action.  Its right side is
exactly the displayed conclusion, simultaneously for all \(u\)
and all maps \(\mathrm{rc}\) satisfying \(\mathrm{IsRecT}\).
We do not replace that right side by a statement for one chosen
correspondence or by a statement only about unramified parameters.
No finite-length or admissibility hypothesis on the product
module is needed by this application.
\end{proof}
% END LTP17

% BEGIN LTP18
\paragraph{LTP18. Every original Teichmuller family uses the same integral type.}
\label{claim:LTP18}
Let \(r\) be the original residual representation and suppose
\(\tau_0\) satisfies its original \(\mathrm{TypeForAlt}(r,\iota)\),
as supplied by ITA12.  For every original family \(t_0\) with
\(r.\mathrm{IsTeich}(t_0)\), define
\[
 \chi_{u,i}(\gamma)
       =\iota\!\left(\operatorname{algebraMap}_{O,\overline{\mathbf Q}_\ell}
                      \bigl(t_{0,i}(V.\mathrm{pr}(\gamma))\bigr)\right),
                         \qquad\gamma\in I_{s.A(u)} .
\]
For the same \(Y\ne0\) and any family of its local components,
LTP17's \(\mathrm{IsTypeWD}\) conclusion holds for these
\(\chi_{u,i}\), at every type place and for every original
\(\mathrm{IsRecT}\) map.
\begin{proof}
This \(\chi\) is the literal pushed-character family in the
definition of \(\mathrm{TypeForAlt}\): the inertia inclusion
into the original absolute Galois group is followed by
\(V.\mathrm{pr}\), then by \(t_{0,i}\), and finally by the
fixed coefficient inclusion and \(\iota\).
The asserted type condition for every such \(t_0\) is a field
of the already supplied \(\mathrm{TypeForAlt}\).
Apply LTP17.  The integral module, \(\tau_0\), \(O\), \(\iota\),
the vector space \(Y\), and the product action \(J\) are unchanged
as \(t_0\) and the local correspondence vary.
In particular the argument retains coefficient prime three,
repeated inertia characters, an empty type set, and every
original allowed rank.  It asserts the Weil--Deligne type
conclusion; descent of a Galois representation and its inertia
conjugating matrices are separate subsequent steps.
\end{proof}
% END LTP18

\subsection{Producing the local components required by the assembly}

% BEGIN LTP19
\paragraph{LTP19. An admissible local component inside the actual automorphic representation.}
\label{claim:LTP19}
For every original automorphic representation \(\sigma\) and
every split finite place \(w\), there is a smooth irreducible
admissible representation \(\tau_w\) of
\(\operatorname{GL}_n(L_w)\) and a nonzero local intertwiner
\(\tau_w\to\sigma\).  This proves the existence and admissibility
clauses of \(\mathrm{DefU.exists\_locComp\_split}\).
The argument uses only the identity-form group and therefore
applies to every CM field and every \(n\ge0\), as in that endpoint.
\begin{proof}
First let \(n\ge1\), and put \(J_w=\operatorname{GL}_n(L_w)\).
LTP2's coordinate construction is a topological product
decomposition
\[
                    G_f=j_w(J_w)\times G_f^{(w)} ,
\]
where the second factor has identity coordinates above
\(w|_{L^+}\).  Both factors commute with \(G_\infty\).

Choose a nonzero vector of \(\sigma\).  Its finite-dimensional
\(G_\infty\)-orbit span contains a nonzero irreducible
\(G_\infty\)-subrepresentation \(S\): the positive invariant form
of ADP1 allows finite-dimensional invariant orthogonal
complements, or one may choose a minimal nonzero invariant
subspace.  Thus the inclusion \(S\hookrightarrow\sigma\) is
an element of
\[
                \mathcal E=\operatorname{Hom}_{G_\infty}(S,\sigma).
\]
Choose a basis of \(S\).  The images of its finitely many basis
vectors have open finite stabilizers by the original automorphic
smoothness condition.  Their intersection contains
\(C_w\times K^{(w)}\), with \(C_w\) compact open in \(J_w\)
and \(K^{(w)}\) compact open in \(G_f^{(w)}\).
Here compact open subgroups are obtained by integral matrices
and principal congruence conditions in the actual finite
coordinates.  Hence
\[
                       M=\mathcal E^{K^{(w)}}\ne0
\]
and \(M\) is stable under \(J_w\).

Every vector of \(M\) is smooth under \(J_w\): intersect the
open finite stabilizers of its values on a basis of \(S\), and
pull back along the continuous \(j_w\).
For every compact open \(C\subset J_w\),
\[
    M^C\ \subset\
       \operatorname{Hom}_{G_\infty}(S,\mathcal A)^{
                                      j_w(C)K^{(w)}} .
\]
The group \(j_w(C)K^{(w)}\) is compact open in the actual \(G_f\).
ADP2 therefore makes the right side finite-dimensional.
This proves the exact local admissibility bound for \(M\);
it does not infer a finite multiplicity statement merely from
compactness of the quotient.

If \(s_1,\ldots,s_e\) is a basis of \(S\), give \(M\) the form
\[
            (T,T')=\sum_{b=1}^e
                         \langle T(s_b),T'(s_b)\rangle_{\mathcal A}.
\]
ADP1 makes it positive definite and \(J_w\)-invariant.
Compact subgroups of \(J_w\) have finite orbits on every smooth
vector of \(M\).  For a compact open \(C\), \(M^C\) is finite
dimensional as just proved.  The algebraic projection argument
PKT5--PKT6 consequently expresses every vector of \(M\) as a
finite sum of vectors in irreducible \(J_w\)-submodules.
All of its hypotheses have been verified here on \(M\).
For clarity, that argument projects inside the finite space
\(M^C\), uses finite compact orbits to make the residual
orthogonal to an entire invariant submodule, and inducts on the
dimension of the fixed space; no Hilbert completion or closed
subrepresentation assumption is used.

Choose a nonzero irreducible submodule \(U\subset M\).
It is smooth, and \(U^C\subset M^C\) is finite-dimensional for
every compact open \(C\), so it is admissible.
There are \(T\in U\) and \(s_0\in S\) with \(T(s_0)\ne0\).
Evaluation
\[
                       U\longrightarrow\sigma,\qquad
                                T\longmapsto T(s_0)
\]
is \(J_w\)-equivariant, since \(J_w\) acts only on the target
of the Hom space.  It is nonzero, so its kernel is zero by
irreducibility of \(U\).  Regard \(U\) with this action as the
smooth irreducible \(\tau_w\).  Since every
\(x\in\mathrm{locU}(n,w)\) equals \(j_w(\mathrm{cLoc}_w x)\),
the evaluation map is precisely a nonzero intertwiner in the
source's \(\mathrm{IsCompAlong}\).

The group, finite-class and multiplicity arguments used here
depend only on a CM field and its identity-form unitary group;
the auxiliary Hida data are absent from their statements and
proofs.  This supplies the asserted general CM-field scope.
For \(n=0\), both adelic and local matrix groups are trivial,
the nonzero irreducible automorphic space is the one-dimensional
space of functions on a point, and its one-dimensional local
representation is smooth and admissible.  This treats the
remaining rank in the literal \texttt{exists\_locComp\_split}
statement.
\end{proof}
% END LTP19

% BEGIN LTP20
\paragraph{LTP20. The exact unramified--spherical equivalence for every local component.}
\label{claim:LTP20}
For every smooth irreducible local component \(\tau\) of
\(\sigma\) at a split place \(w\),
\[
            \mathrm{IsUnramU}(\sigma,w)
                 \quad\Longleftrightarrow\quad
                   \tau^{KGL_w}\ne0 .
\]
Together with LTP19 this proves the full original
\(\mathrm{DefU.exists\_locComp\_split}\), including admissibility.
\begin{proof}
If \(\mathrm{IsUnramU}\) holds, its nonzero function is fixed
by \(j_w(KGL_w)\).  LTP3 makes \(\sigma\) restricted along
\(j_w\) \(\tau\)-isotypic.  The equality \(KGL_w=KO_w\) and
compactness from LTP5 therefore imply \(\tau^{KGL_w}\ne0\).

Conversely, let \(0\ne v\in\tau^{KGL_w}\).
The supplied component map \(f:\tau\to\sigma\) is nonzero.
Its kernel is invariant under the local group, so irreducibility
makes \(f\) injective.  Thus \(\phi=f(v)\ne0\).
For every \(x\in\mathrm{locU}(n,w)\) whose \(w\)-coordinate
belongs to \(KGL_w\), equivariance gives
\[
       \mathscr R_x\phi
                   =f\bigl(\tau(\mathrm{cLoc}_w x)v\bigr)
                   =f(v).
\]
This is exactly the witness demanded by \(\mathrm{IsUnramU}\).
The equivalence holds for every supplied component, not only
for the admissible component chosen in LTP19.
\end{proof}
% END LTP20

% BEGIN LTP21
\paragraph{LTP21. One family of components supplies all the original local fields.}
\label{claim:LTP21}
Fix \(Y\ne0\).  One may choose a single family
\((\tau_w)_{w\text{ split}}\) before using any of the following
local conclusions.  Each \(\tau_w\) is smooth irreducible and
admissible and is a component of this same \(\sigma\).  For this
one family the source's local requirements hold simultaneously:
\[
\begin{array}{c|l}
 \text{place or condition}&\text{conclusion on the original objects}\\ \hline
 w=s.w(v),\ v\mid\ell&\mathrm{IsoSum}(H.\mathrm{atV}(v),Y,\tau_w)\\
 w=Q.wq(q)&\mathrm{IsoSum}(H.\mathrm{atW}(w),Y,\tau_w)\\
 w|_{L^+}\notin V.S',\ w\text{ split}
       &\tau_w^{KGL_w}\ne0,\ \text{and the LTP11 spherical transfer}\\
 a\in I&\mathrm{HasWeightAt}(\sigma,\iota\circ a.\mathrm{res},
                                               \lambda(a.\mathrm{res})^\vee)\\
 w=s.w(v_0)&\tau_w^{K_c(L_w,n,1)}\ne0\\
 u\in s.St&\text{LTP17 for every original }\mathrm{IsRecT}
                 \text{ and every supplied type family.}
\end{array}
\]
In addition \(\mathrm{IsUnramU}(\sigma,w)\) holds at every
finite \(w\) outside \(V.S'\), whether or not \(w\) splits.
\begin{proof}
Apply choice once to LTP19, indexed by the set of split places.
Retain the resulting local representations and component maps
throughout.  FNL4 identifies the selected \(\ell\)-places,
auxiliary places, type places and \(v_0\) as members of this
index set.  LTP7, LTP9, LTP12, LTP14 and LTP15 give the
corresponding statements, and LTP20 changes LTP12 into spherical
fixedness at a split outside place.  LTP11 transfers every
given spherical scalar, on the original operator and every
coefficient vector in \(Y\).  Restrict the same family to the
type places in LTP17; LTP18 makes the conclusion universal over
the original Teichmuller families when the supplied integral
type is used.  FNL16 supplies the original \(H.\mathrm{Ax}\)
needed by LTP9 and LTP11.

This is one collection of local components of one original
automorphic representation, at the same \(E,\xi,\iota,\tau_0\)
and \(Y\).  No independent replacement of components in the
different rows is necessary.  It supplies the local-transport
inputs to the subsequent \texttt{PartU} construction; that
construction still has to provide its global Galois representation,
local--global and Frobenius comparisons, quotient descent,
pairing, and point-property record.
\end{proof}
% END LTP21
